\begin{ouverture}[college4]{18}{Signaux et information}{Partie 4 · Des signaux pour observer et communiquer}{Comment un signal devient-il un message ?}
\objectif{Identifier émetteur, support de propagation et récepteur.}
\objectif{Distinguer information et grandeur physique du signal.}
\objectif{Lire un codage et discuter la qualité d’une mesure.}
\prerequis{Sources lumineuses et sonores, graphique et proportionnalité.}
\end{ouverture}
\pcPage{Le cours}{Signaux et information}
\phantomsection\label{ess-18}
\essSection{1}{Une chaîne de communication}
Un émetteur produit un signal ; un récepteur le détecte et permet d’en extraire une information grâce à un code ou un traitement. Un signal peut être lumineux, sonore, électrique ou radio. Les ondes radio appartiennent à la famille des rayonnements électromagnétiques : elles ne sont pas des sons et peuvent se propager dans le vide. Le signal est un phénomène physique ; l’information est le sens ou la donnée qu’on lui associe.\par
\pcFigure{figures/c20-chaine.pdf}{Une communication peut convertir plusieurs fois la nature du signal entre la source du message et son destinataire.}
\essSection{2}{Coder un message}
Un code établit une correspondance connue entre des signaux et des symboles. Deux états distincts permettent de coder 0 et 1 ; une suite de bits peut ensuite représenter plusieurs symboles. Un signal n’a pas de sens unique sans convention : un voyant rouge peut annoncer une alerte ou un état choisi par le concepteur. Le codage doit préciser les durées, l’ordre des symboles et leur séparation.\par
\clearpage
\essSection{3}{Mesurer avec un signal}
Un capteur transforme une grandeur physique en une sortie exploitable. Un étalonnage relie des valeurs connues de la grandeur aux valeurs du signal. Dans une plage donnée, une relation peut être proportionnelle ou affine, mais ce n’est pas automatique. Un signal analogique varie continûment dans son modèle ; un traitement numérique utilise des valeurs discrètes. Cette distinction ne signifie pas que le numérique est toujours exact.\par
\pcFigure{figures/c20-signaux.pdf}{À gauche : variation continue. À droite : états de code 0 ou 1 ; les valeurs et le sens du message dépendent de la convention.}
\essSection{4}{Détecter les limites}
Bruit, obstacle, distance, mauvaise orientation et perturbations peuvent dégrader la réception. Répéter un message ou ajouter une vérification aide parfois à repérer une erreur, sans la supprimer toujours. Pour une mesure, on contrôle étalonnage, plage, résolution et temps de réponse. Une donnée affichée avec beaucoup de chiffres ne garantit pas une grande exactitude. Expliquer la limite du dispositif fait partie d’une conclusion scientifique.\par
\begin{pcEncadre}{method}{Un exemple expliqué}Un capteur scolaire donne $U=0{,}020\,T+0{,}50$, avec U en V et T en degrés Celsius, entre 0 et 50 °C. Pour $U=1{,}10\unite V$, $T=(1{,}10-0{,}50)/0{,}020=30\,^{\circ}\mathrm C$. Ce calcul utilise l’étalonnage fourni ; il ne permet pas d’extrapoler à une température arbitraire.\end{pcEncadre}
\pcPage{Expérimenter et retenir}{Signaux et information}
\phantomsection\label{activite-18}
\begin{pcEncadre}{method}{Investigation en classe}\textbf{Mission.} Inventer un code lumineux à deux états pour quatre messages, puis le transmettre à faible distance avec une petite lampe adaptée. Échanger le code avant la transmission. Tester un message, puis introduire sur document une erreur de symbole. Proposer un moyen de détection et discuter sa limite. \textbf{À rendre :} table du code, chronologie des états, chaîne émetteur-récepteur et bilan de fiabilité. Aucun émetteur radio n’est fabriqué et aucune lumière intense n’est utilisée.\end{pcEncadre}
\pcSousTitre{Ce que je dois savoir expliquer}\begin{questions}
\item Le signal est un phénomène physique portant une information codée.
\item Un récepteur doit disposer d’un code ou d’un traitement.
\item Les ondes radio ne sont pas des ondes sonores.
\item Une mesure par capteur exige un étalonnage et une plage d’utilisation.
\item Les perturbations peuvent modifier le signal et l’information reçue.
\end{questions}
\begin{pcEncadre}{attention}{Une erreur à éviter}Un affichage numérique de 23,456 °C ne prouve pas que la température est connue au millième de degré : la qualité du capteur et de l’étalonnage limite le résultat.\end{pcEncadre}
\pcRepere{cor-activite-18}{Pistes pour l’investigation}\pcGrille
\clearpage\markright{Je démarre}
\rubrique{Je démarre}
\par\noindent\begin{minipage}{\linewidth}
\begin{exercice}[1]{Les rôles}
\begin{questions}
\item Dans une sonnette, identifier un émetteur sonore et un récepteur humain.
\item Quelle information peut être transmise ?
\end{questions}
\end{exercice}
\end{minipage}
\par\noindent\begin{minipage}{\linewidth}
\begin{exercice}[1]{Radio et son}
\begin{questions}
\item Une onde radio est-elle un son très aigu ?
\item Peut-elle traverser le vide ?
\end{questions}
\end{exercice}
\end{minipage}
\par\noindent\begin{minipage}{\linewidth}
\begin{exercice}[1]{Deux états}
\begin{questions}
\item Que signifie un code à deux états ?
\item Le voyant allumé signifie-t-il partout « danger » ?
\end{questions}
\end{exercice}
\end{minipage}
\par\noindent\begin{minipage}{\linewidth}
\begin{exercice}[1]{Le capteur}
\begin{questions}
\item Quel est le rôle d’un capteur ?
\item Pourquoi l’étalonner ?
\end{questions}
\end{exercice}
\end{minipage}
\par\noindent\begin{minipage}{\linewidth}
\begin{exercice}[1]{Une perturbation}
\begin{questions}
\item Nommer une cause de mauvaise réception lumineuse.
\item Et une cause de mauvaise réception sonore ?
\end{questions}
\end{exercice}
\end{minipage}
\clearpage\markright{Je m’entraîne}
\rubrique{Je m’entraîne}
\par\noindent\begin{minipage}{\linewidth}
\begin{exercice}[2]{Lire un code}
\begin{questions}
\item Le code fourni est A = 00, B = 01, C = 10, D = 11. Décoder 01 10 00.
\item Pourquoi les espaces sont-ils utiles ?
\end{questions}
\end{exercice}
\end{minipage}
\par\noindent\begin{minipage}{\linewidth}
\begin{exercice}[2]{Compter des messages}
\begin{questions}
\item Combien de suites différentes peut-on faire avec deux bits ?
\item Donner les suites.
\end{questions}
\end{exercice}
\end{minipage}
\par\noindent\begin{minipage}{\linewidth}
\begin{exercice}[2]{Un étalonnage simple}
\begin{questions}
\item Un capteur donne $U=0{,}10\,x$ avec U en V et x dans son unité. Calculer U pour x = 12.
\item Retrouver x pour U = 0,8 V.
\end{questions}
\end{exercice}
\end{minipage}
\par\noindent\begin{minipage}{\linewidth}
\begin{exercice}[2]{Une sortie affine}
\pcFigureEx{figures/ex-etalonnage.pdf}{}
\begin{questions}
\item Pour $U=0{,}020T+0{,}50$, calculer U à 20 °C.
\item La relation est-elle proportionnelle entre U et T ?
\end{questions}
\end{exercice}
\end{minipage}
\par\noindent\begin{minipage}{\linewidth}
\begin{exercice}[2]{Une valeur reçue}
\begin{questions}
\item Avec l’étalonnage précédent, retrouver T pour $U=1{,}30\unite V$.
\item La valeur est-elle dans la plage 0 à 50 °C ?
\end{questions}
\end{exercice}
\end{minipage}
\clearpage\markright{Je raisonne}
\rubrique{Je raisonne}
\par\noindent\begin{minipage}{\linewidth}
\begin{exercice}[3]{Une erreur de bit}
\begin{questions}
\item Le code de l’exercice 6 est utilisé. B = 01 est reçu 11. Quel message est lu ?
\item Le récepteur peut-il détecter l’erreur avec cette seule table ?
\end{questions}
\end{exercice}
\end{minipage}
\par\noindent\begin{minipage}{\linewidth}
\begin{exercice}[3]{Répéter et comparer}
\begin{questions}
\item Un message est reçu deux fois : 01 10 puis 01 11. Que montre la comparaison ?
\item Permet-elle de savoir quelle version est correcte ?
\end{questions}
\end{exercice}
\end{minipage}
\par\noindent\begin{minipage}{\linewidth}
\begin{exercice}[3]{Hors de la plage}
\pcFigureEx{figures/ex-hors-plage.pdf}{La zone colorée correspond au seul domaine étalonné.}
\begin{questions}
\item Un élève utilise $U=0{,}020T+0{,}50$ pour U = 3 V, alors que l’étalonnage vaut de 0 à 50 °C. Quel calcul obtient-il ?
\item Pourquoi ne doit-il pas annoncer ce résultat comme une mesure fiable ?
\end{questions}
\end{exercice}
\end{minipage}
\par\noindent\begin{minipage}{\linewidth}
\begin{exercice}[3]{Une résolution fournie}
\begin{questions}
\item Un capteur affiche la température par pas de $0{,}5\,^{\circ}\mathrm C$. Peut-il distinguer à l’affichage 20,1 et 20,2 °C ?
\item Cette résolution résume-t-elle toute la qualité de mesure ?
\end{questions}
\end{exercice}
\end{minipage}
\par\noindent\begin{minipage}{\linewidth}
\begin{exercice}[3]{Analogique ou numérique}
\begin{questions}
\item Une tension varie continûment puis un convertisseur donne des nombres arrondis. Décrire les deux représentations.
\item Les arrondis constituent-ils une erreur humaine ?
\end{questions}
\end{exercice}
\end{minipage}
\pcPage{Défi du chapitre}{Signaux et information}
\phantomsection\label{defi-ex-18}
\pcDefiTitre{18}{Le protocole du phare}
\begin{questions}
\item On code quatre messages par deux bits, avec 1 = lampe allumée et 0 = éteinte. Chaque bit dure 1 s. Combien de temps dure un message, hors séparation ?
\item Proposer un signal de début pour éviter de confondre silence et 00.
\item Pourquoi la seule répétition ne garantit-elle pas une correction parfaite ?
\end{questions}
\pcGrille
